% Esperienza professionale — standalone section.
% Date column is as wide as the longest date and right-aligned.
% Load cabin before cochineal. The other order segfaults this Tectonic.
\documentclass[10pt,a4paper]{article}

\usepackage[margin=0.9cm]{geometry}
\usepackage[italian]{babel}
\usepackage{cabin}
\usepackage{cochineal}
\usepackage{fontawesome5}
\usepackage{xcolor}
\usepackage{enumitem}
\usepackage{tabularx}

\definecolor{accent}{rgb}{0.57098,0.02196,0.19766}
\definecolor{rulecolor}{rgb}{0.72,0.8047,0.42589}

\pagestyle{empty}
\setlength{\parindent}{0pt}
\setlength{\parskip}{0pt}

\setlist[itemize]{%
  leftmargin=2.5em,
  labelsep=0.5em,
  label=\textbullet,
  topsep=5pt,
  itemsep=4pt,
  parsep=0pt,
  partopsep=0pt,
  before=\footnotesize,
  font=\footnotesize,
}

\newcommand{\sep}{\textcolor{accent}{\small\faBookmark}}
\newcommand{\role}[1]{\textbf{#1}}
\newcommand{\org}[1]{\textit{#1}}

\begin{document}

\noindent{\sffamily\bfseries\Large Esperienza professionale}\par
% Bar bottom is 4.33pt under the title baseline; 9.6pt then meets the date row.
\vskip-\prevdepth
\vskip-4.185pt
\noindent\textcolor{rulecolor}{\rule{\textwidth}{2.5pt}}\par
\vskip 9.6pt

\noindent
\begin{tabularx}{\textwidth}{@{}>{\raggedleft\arraybackslash}r@{\hspace{0.75cm}}l@{\hspace{9.46pt}}X@{}}
Gennaio 2025 -- in corso
  & \sep
  & \role{Data Scientist / Data Engineer, \org{Complexa Labs}, Bari (IT)}
    \begin{itemize}
      \item Realizzazione di soluzioni di analytics come consulente freelance per clienti enterprise (fatturato superiore a 1 milione di euro).
      \item Sviluppo di un motore ATS di time-to-fill: modello di sopravvivenza a tempo discreto score-driven (GAS) su un pannello posizione vacante--giorno, con posizioni aperte censurate a destra e covariate giornaliere del funnel di selezione, in Polars, messo in produzione su SageMaker.
      \item Reingegnerizzazione delle pipeline in operazioni native Spark (PySpark): join nativi, funzioni finestra e broadcast hint, per ottimizzare le prestazioni.
      \item Sviluppo di un simulatore di scenari macroeconomici, con orchestrazione distribuita delle pipeline, per prevedere gli impatti sul business al variare delle condizioni di mercato.
      \item Sviluppo di un modello di Customer Lifetime Value (CLTV) che combina analisi di sopravvivenza mixture-cure, regressione del margine con gradient boosting e segmentazione calibrata, per la previsione dei ricavi e il targeting del Next Best Product.
      \item Sviluppo di algoritmi NLP per l'automazione dell'elaborazione documentale, messi in produzione come microservizi FastAPI in Python.
    % Keep 0.5pt after dropping the list's trailing skip (endpbox would unskip it).
    \end{itemize}\unskip\vskip0.5pt\penalty0\tabularnewline
Gennaio 2024 -- novembre 2024
  & \sep
  & \role{Data Scientist, \org{acmetric B.V.}, Amsterdam (NL)}
    \begin{itemize}
      \item Sviluppo e manutenzione di un motore di pricing e previsione della domanda per Adidas International B.V., basato su modelli di Causal AI.
      \item Progettazione e messa in produzione di un modello di previsione di churn e contenzioso per LATAM Airlines, con Gradient Boosting (GBM).
    % Keep 0.5pt after dropping the list's trailing skip (endpbox would unskip it).
    \end{itemize}\unskip\vskip0.5pt\penalty0\tabularnewline
Novembre 2022 -- dicembre 2023
  & \sep
  & \role{Data Scientist, \org{MOOS Solutions B.V.}, Amsterdam (NL)}
    \begin{itemize}
      \item Sviluppo e manutenzione di un motore di intelligenza artificiale per analisi di inventario in tempo reale, basato su tecniche Mixture of Experts (MoE).
    \end{itemize}
\end{tabularx}

\end{document}
